% Ad Atticum 16.5 (ad-atticum-16-05) --- English translation
% Translated by Claude (Anthropic) from the Latin (Perseus).
% Style guide: STYLE.md.

\ciceroLetterOpener{CICERO ATTICO salutem}

\ciceroSection{1} Brutus by now was expecting a letter from you. I had brought him no new word about Accius's \textit{Tereus}; he had supposed it was Brutus's own. Still, some rumour had reached him --- I do not know which --- that at the opening performance the Greeks had not turned out in numbers; which did not surprise me at all, for you know what I think of Greek games.

\ciceroSection{2} Now hear what counts more than everything else. Quintus spent several days with me --- and would have stayed more, had I wished it; but for the time he was with me, it is incredible how he pleased me in every respect, and most of all in the very respect where he used to satisfy me least. He is so wholly changed --- partly by certain writings of mine I had to hand, partly by my continual conversation and admonition --- that he will face the state with the spirit we could wish him to have. When he had not merely assured me of this but won me over to believe it, he begged me at length and in earnest to give you my pledge that he would prove worthy both of you and of us; not that he was asking you to trust him at once, but that, when you had looked him over for yourself, you should then love him. Unless he had won my confidence, and unless I judged that what I am saying would hold firm, I should not have done what I am about to tell you. For I took the young man with me to Brutus. So thoroughly did he prove what I am here writing to you, that Brutus took it on faith himself, would not accept me as guarantor, and praising the lad in the warmest terms made affectionate mention of you, embraced him, and dismissed him with a kiss. For all that, though there is more for me to congratulate you on than to ask of you, I do still ask this: that, whatever lapses in him before now seemed to come of the want of steadiness that goes with his age, you should judge he has shed; and trust me that your authority will count for much --- or rather for the most --- in confirming his fixed purpose.

\ciceroSection{3} Though I have often dropped to Brutus the suggestion of a convoy \textit{[Greek: homoploiai]}, he did not, as I had supposed he would, seem to leap at it. I judged that he was somewhat up in the air \textit{[Greek: meteoroteron]} --- and by Hercules he was, and chiefly because of the games. But when I had got back to the villa, Gnaeus Lucceius, who is much in Brutus's company, told me that he is much delayed --- not from any wish to evade, but from waiting in case something turns up. So I am hesitating whether to make for Venusia and there wait for news of the legions; or, if they are away (as some think), then for Hydruntum; and if neither is safe \textit{[Greek: asphales]}, to come back here. You think I am joking? I will die if anyone but you holds me here. Look around you --- only do it before I blush.

\ciceroSection{4} O days under Lepidus's auspices, prettily laid out and aptly suited to a plan for our coming back! Your letter gives a great lurch \textit{[Greek: rhope]} toward setting out. If only you yourself were there! But do what you judge expedient.

\ciceroSection{5} I am waiting for Nepos's letter. Is the man so fond of my work? --- when he does not think what I take most pride in \textit{[Greek: gauriao]} worth reading. And you say ``next after the matchless one'' \textit{[Greek: met' amymona]} --- you are matchless \textit{[Greek: amymon]}, but he is divine \textit{[Greek: ambrotos]}. Of my letters there is no collection \textit{[Greek: synagoge]} --- but Tiro has a thing like seventy, and indeed some must be got from you. These I must look over and correct; only then will they be published.
